\part{Die Gestalt}

\chapter{Wie groß eine Varietät ist}

\section{Monome zählen}

\begin{merksatz}
\boxtitel{Die Leitfall-Kurve, gezählt}
Die Gröbner-Basis des Leitfalls in grevlex hat die Leitmonome $k_{10}$, $k_{21}$ und $k_{12}k_{20}$. Man zähle die Monome vom Gesamtgrad $\le d$, die durch \kursiv{keines} davon teilbar sind --- die Standardmonome:

\begin{center}
{\small
\renewcommand{\arraystretch}{1.46}
\begin{tabular}{@{}P{3.0cm}P{0.7cm}P{0.7cm}P{0.7cm}P{0.7cm}P{0.7cm}P{0.7cm}P{1.2cm}@{}}
\toprule
$d$ & $0$ & $1$ & $2$ & $3$ & $4$ & $5$ & \textbf{Formel}\\
\midrule
Leitfall-Kurve & $1$ & $3$ & $5$ & $7$ & $9$ & $11$ & $2d + 1$\\
verdrehte Kubik & $1$ & $4$ & $7$ & $10$ & $13$ & $16$ & $3d + 1$\\
Kubik $y^{2} = x^{3} - 2x$ & $1$ & $3$ & $6$ & $9$ & $12$ & $15$ & $3d$\\
Sattelfläche & $1$ & $4$ & $9$ & $16$ & $25$ & $36$ & $(d+1)^{2}$\\
Ebene $\cup$ Gerade & $1$ & $4$ & $8$ & $13$ & $19$ & $26$ & $\tfrac12d^{2} + \tfrac52d + 1$\\
zwei Punkte (Aufg.~1.3) & $1$ & $2$ & $2$ & $2$ & $2$ & $2$ & $2$\\
Kreis $\cap$ Parabel & $1$ & $3$ & $4$ & $4$ & $4$ & $4$ & $4$\\
\bottomrule
\end{tabular}}
\end{center}

\begin{center}
\begin{minipage}{0.9\textwidth}
\textbf{Für große $d$ wird die Zählung ein Polynom in $d$. Sein Grad ist die Dimension der Varietät --- $1$ für Kurven, $2$ für Flächen, $0$ für Punktmengen. Sein Leitkoeffizient, mal Fakultät des Grades, ist der Grad der Varietät: $2$ für die Hyperbel, $3$ für die Kubiken, $4$ für die vier Schnittpunkte.}
\end{minipage}
\end{center}

\begin{center}
\begin{minipage}{0.88\textwidth}
\kursiv{\textbf{Dimension ist ein Wort, das jeder versteht und niemand definieren kann, ohne zu rechnen.} Dieses Kapitel gibt drei Definitionen --- über das Wachstum der Funktionen auf der Varietät (Hilbert-Polynom), über Ketten von Teilvarietäten (Krull) und über freie Veränderliche (Transzendenzgrad) --- und beweist, dass die erste berechenbar ist und mit den anderen übereinstimmt (bis auf eine Auslassung). Die Zahl $2d + 1$ für den Leitfall ist der Satz \enquote{nicht identifizierbar, mit einem Freiheitsgrad} als Polynom; und ihr Leitkoeffizient $2$ sagt, dass eine allgemeine zusätzliche Messung die Zahl der passenden Modelle auf zwei reduzieren würde --- nicht auf eins.}
\end{minipage}
\end{center}
\end{merksatz}

\section{Die Hilbert-Funktion}

\begin{definition}
\boxtitelm{Definition 7.1 --- Hilbert-Funktion und Hilbert-Polynom}
Für ein Ideal $I\subseteq k[x_1, \ldots, x_n]$ sei $k[x]_{\le d}$ der Raum der Polynome vom Gesamtgrad $\le d$ und $I_{\le d} = I\cap k[x]_{\le d}$. Die \Begriff{(affine) Hilbert-Funktion} ist
\[ {}^{a}\!H_I(d) = \dim_k k[x]_{\le d}/I_{\le d} = \dim_k k[x]_{\le d} - \dim_k I_{\le d} . \]
Die Hilbert-Funktion einer Varietät $V$ ist die von $I(V)$: die Zahl der linear unabhängigen Polynomfunktionen vom Grad $\le d$ auf $V$.
\end{definition}

\begin{merksatz}
\boxtitel{Satz 7.2 --- Macaulay: Standardmonome zählen}
Sei $>$ eine gradkompatible Monomordnung (grlex oder grevlex: höherer Gesamtgrad ist größer). Dann ist
\[ {}^{a}\!H_I(d) = \#\{\,\text{Monome } x^{\alpha} : |\alpha|\le d,\ x^{\alpha}\notin\mathrm{LM}(I)\,\} , \]
und die Standardmonome --- die nicht in $\mathrm{LM}(I)$ --- bilden eine Basis von $k[x]/I$ als $k$-Vektorraum.

\schritt{Beweis}
Die Monome vom Grad $\le d$ in $\mathrm{LM}(I)$ sind genau die Leitmonome der Elemente von $I_{\le d}$: Ist $f\in I_{\le d}$, so hat $\mathrm{LM}(f)$ Grad $\le d$ (gradkompatibel!); ist umgekehrt $x^{\alpha}\in\mathrm{LM}(I)$ mit $|\alpha|\le d$, so ist $x^{\alpha} = \mathrm{LM}(f)$ für ein $f\in I$, und da die Ordnung gradkompatibel ist, hat $f$ Grad $|\alpha|\le d$. Nun ist $\dim I_{\le d}$ gleich der Zahl der verschiedenen Leitmonome in $I_{\le d}$: Die Elemente von $I_{\le d}$ mit verschiedenen Leitmonomen sind linear unabhängig, und durch Gauß-Elimination nach Leitmonomen erhält man eine Basis mit lauter verschiedenen Leitmonomen. Also $\dim I_{\le d} = \#\{x^{\alpha}\in\mathrm{LM}(I) : |\alpha|\le d\}$, und die Formel folgt. --- Basis: Jedes $f$ hat einen eindeutigen Rest modulo einer Gröbner-Basis (Satz 5.2), und der ist eine Kombination von Standardmonomen; Standardmonome sind linear unabhängig modulo $I$, weil eine Kombination in $I$ ein Leitmonom in $\mathrm{LM}(I)$ hätte. $\blacksquare$

\schritt{Warum gradkompatibel}
In lex kann das Leitmonom kleineren Grad haben als ein anderer Term: $\mathrm{LM}(x + y^{5}) = x$ für $x > y$. Dann liegt $x$ in $\mathrm{LM}(I)$, obwohl kein Element vom Grad $1$ in $I$ liegt, und die Zählung geht schief. Mit grevlex hat das Leitmonom stets den Gesamtgrad des Polynoms, und die Formel stimmt; deshalb ist grevlex die Ordnung für Dimension und Grad.

\begin{center}
\begin{minipage}{0.88\textwidth}
\kursiv{\textbf{Der Satz macht die Hilbert-Funktion zu einer Abzählaufgabe über ein Monomideal} --- und Monomideale kann man zählen: Für $\mathrm{LM}(I) = \langle k_{10}, k_{21}, k_{12}k_{20}\rangle$ sind die Standardmonome $k_{12}^{a}$ und $k_{20}^{b}$ ($a, b\ge0$, nicht beide $>0$); vom Grad $\le d$ gibt es $1 + d + d = 2d + 1$ davon. \checkmark\ Für die verdrehte Kubik mit Leitmonomen $x^{2}, xy, xz, y^{3}$ sind die Standardmonome $x$ und $y^{a}z^{b}$ mit $a\le2$; Aufgabe 7.1 zählt sie: $3d + 1$.}
\end{minipage}
\end{center}
\end{merksatz}

\begin{merksatz}
\boxtitel{Satz 7.3 --- das Hilbert-Polynom}
Für jedes Ideal $I$ gibt es ein Polynom ${}^{a}\!HP_I(d)$ mit rationalen Koeffizienten und ein $d_0$, so dass ${}^{a}\!H_I(d) = {}^{a}\!HP_I(d)$ für alle $d\ge d_0$. Sein Grad hängt nur von $V(I)$ ab (nicht von $I$), und man definiert
\[ \dim V := \deg{}^{a}\!HP_{I(V)}, \qquad \deg V := (\dim V)!\cdot(\text{Leitkoeffizient von }{}^{a}\!HP_{I(V)}) . \]

\schritt{Beweis für Monomideale}
Sei $M = \langle x^{\alpha(1)}, \ldots, x^{\alpha(r)}\rangle$. Die Monome außerhalb von $M$ zerfallen in endlich viele \enquote{Koordinatenunterräume mit Verschiebung}: Für jede Teilmenge $S\subseteq\{1, \ldots, n\}$ und jeden Exponentenvektor $\beta$ mit $\beta_i = 0$ für $i\in S$ betrachte man die Monome $x^{\beta}\prod_{i\in S}x_i^{a_i}$, $a_i\ge0$. Die Komplementmenge von $M$ ist eine endliche disjunkte Vereinigung solcher Mengen (Induktion nach $n$ und nach der Zahl der Erzeuger; Auslassung 5, \zit{5}, Kap.~9.2). Die Zahl der Monome vom Grad $\le d$ in einer solchen Menge mit $|S| = s$ ist $\binom{d - |\beta| + s}{s}$, ein Polynom vom Grad $s$ in $d$ für $d\ge|\beta|$. Die Summe endlich vieler solcher Polynome ist ein Polynom. $\blacksquare$

\schritt{Für beliebige Ideale}
Nach Satz 7.2 ist ${}^{a}\!H_I = {}^{a}\!H_{\mathrm{LM}(I)}$ für gradkompatible Ordnungen: dieselbe Funktion wie ein Monomideal. $\blacksquare$

\schritt{Dimension hängt nur von $V$ ab}
${}^{a}\!H_{\sqrt I}\le{}^{a}\!H_I$ (größeres Ideal, kleinerer Quotient), und andererseits $I^{m}\subseteq\sqrt I^{\,m}\subseteq I$ für ein $m$ (Basissatz), was ${}^{a}\!H_I(d)\le{}^{a}\!H_{\sqrt I}(md)$ gibt --- beide Hilbert-Polynome haben denselben Grad. Und $V(I) = V(\sqrt I)$, $\sqrt I = I(V)$ über $\mathbb C$. $\blacksquare$

\begin{center}
\begin{minipage}{0.9\textwidth}
\textbf{Dimension misst, wie schnell die Zahl der Funktionen auf $V$ mit dem Grad wächst --- wie $d^{\dim V}$. Auf einer Kurve gibt es linear viele Funktionen vom Grad $\le d$, auf einer Fläche quadratisch viele, auf $n$ Punkten für großes $d$ genau $n$: Mehr als $n$ Werte kann eine Funktion auf $n$ Punkten nicht annehmen.}
\end{minipage}
\end{center}

\begin{center}
\begin{minipage}{0.88\textwidth}
\kursiv{\textbf{Der Grad einer nulldimensionalen Varietät ist die Zahl ihrer Punkte} (mit Vielfachheit, wenn $I$ nicht radikal ist): Für Kreis $\cap$ Parabel ist das Hilbert-Polynom $4$, für die beiden Punkte aus Aufgabe 1.3 ist es $2$ --- das ist der Satz hinter der Regel \enquote{Zahl der Lösungen = Zahl der Standardmonome} aus Aufgabe 5.3. Der Grad einer Kurve ist, wie Kapitel 9 zeigt, die Zahl ihrer Schnittpunkte mit einer allgemeinen Hyperebene; die Hilbert-Funktion zählt ihn, ohne einen Schnitt zu bilden.}
\end{minipage}
\end{center}
\end{merksatz}

\section{Zwei weitere Dimensionen}

\begin{merksatz}
\boxtitel{Satz 7.4 --- drei Definitionen, eine Zahl}
Für eine irreduzible affine Varietät $V$ über $\mathbb C$ stimmen überein:
\begin{center}
{\small
\renewcommand{\arraystretch}{1.5}
\begin{tabular}{@{}P{2.6cm}P{6.2cm}@{}}
\toprule
\textbf{Hilbert} & Grad des Hilbert-Polynoms von $I(V)$\\
\textbf{Krull} & die größte Länge $r$ einer Kette $V_0\subsetneq V_1\subsetneq\cdots\subsetneq V_r = V$ irreduzibler Teilvarietäten (gleichwertig: Ketten von Primidealen in $k[V]$)\\
\textbf{Transzendenz} & die größte Zahl von Koordinatenfunktionen $x_{i_1}, \ldots, x_{i_r}$, zwischen denen auf $V$ keine polynomiale Gleichung gilt ($I(V)\cap k[x_{i_1}, \ldots, x_{i_r}] = \{0\}$)\\
\bottomrule
\end{tabular}}
\end{center}
Der Beweis der Gleichheit ist Auslassung 6 (\zit{5}, Kap.~9.4--9.5; \zit{10}); hier die Richtung, die man braucht.

\schritt{Transzendenz $\le$ Hilbert}
Gilt zwischen $x_{i_1}, \ldots, x_{i_r}$ keine Gleichung, so sind alle Monome in ihnen linear unabhängig modulo $I(V)$ --- das sind $\binom{d + r}{r}$ vom Grad $\le d$, ein Polynom vom Grad $r$. Also wächst ${}^{a}\!H$ mindestens wie $d^{r}$: $\deg{}^{a}\!HP\ge r$. $\blacksquare$

\schritt{Hilbert $\le$ Transzendenz, Idee}
Die Standardmonome liegen in endlich vielen verschobenen Koordinatenunterräumen (Beweis von Satz 7.3); der größte hat $s = \deg{}^{a}\!HP$ Veränderliche $x_{i_1}, \ldots, x_{i_s}$, und für diese gilt: Kein Monom in ihnen allein liegt in $\mathrm{LM}(I)$, also kein Polynom in ihnen allein in $I$ (sein Leitmonom wäre eines). Also $I\cap k[x_{i_1}, \ldots, x_{i_s}] = \{0\}$. $\blacksquare$

\schritt{Leitfall}
Freie Veränderliche: $k_{12}$ allein (oder $k_{20}$ allein) --- zwischen $k_{12}$ und $k_{20}$ gilt $g = 0$, also höchstens eine. Dimension $1$, dreimal. \checkmark

\begin{center}
\begin{minipage}{0.88\textwidth}
\kursiv{\textbf{Die dritte Definition ist die, die der Praktiker benutzt}: Dimension = Zahl der frei wählbaren Parameter, und die freien Parameter liest man an der Gröbner-Basis ab --- es sind die Veränderlichen, die in keinem Leitmonom vorkommen, in der größten solchen Menge. Die zweite ist die, die der Algebraiker benutzt, weil sie für beliebige Ringe Sinn hat. Die erste ist die, die der Rechner benutzt. Dass alle drei dieselbe Zahl geben, ist ein Satz, kein Zufall, und er ist das Fundament, auf dem \enquote{Dimension} ein Begriff und nicht eine Anschauung ist.}
\end{minipage}
\end{center}
\end{merksatz}

\section{Fasern}

\begin{merksatz}
\boxtitel{Satz 7.5 --- die Dimension der Fasern}
Sei $\varphi\colon V\to W$ eine polynomiale Abbildung irreduzibler Varietäten mit dichtem Bild ($\overline{\varphi(V)} = W$). Dann gilt für jeden Punkt $w$ im Bild und jede Komponente $F$ der Faser $\varphi^{-1}(w)$:
\[ \dim F\ge\dim V - \dim W, \]
und für $w$ in einer nichtleeren offenen Teilmenge von $W$ gilt Gleichheit (Auslassung 7, \zit{10}, Kap.~I.6).

\schritt{Der Leitfall}
$\varphi\colon(k_{10}, k_{12}, k_{21}, k_{20})\mapsto(s_1, s_2, s_3)$, $V = \mathbb C^{4}$, $W = \mathbb C^{3}$ (Kapitel 6: keine Gleichung zwischen den $s_i$, also dichtes Bild). Allgemeine Faser: Dimension $4 - 3 = 1$. \checkmark\ Für das Modell mit $k_{20} = 0$: $\mathbb C^{3}\to\mathbb C^{3}$, allgemeine Faser Dimension $0$ --- endlich viele Punkte, und die Hilbert-Funktion sagt, wie viele: einer (Aufgabe 1.2), oder zwei bei Messung im Gewebe (Aufgabe 1.3).

\begin{center}
{\small
\renewcommand{\arraystretch}{1.5}
\begin{tabular}{@{}P{3.2cm}P{1.6cm}P{1.6cm}P{2.0cm}@{}}
\toprule
\textbf{Modell} & $\dim V$ & $\dim W$ & \textbf{Faser}\\
\midrule
vier Parameter, Blut & $4$ & $3$ & Kurve\\
$k_{20} = 0$, Blut & $3$ & $3$ & $1$ Punkt\\
$k_{20} = 0$, Gewebe & $3$ & $3$ & $2$ Punkte\\
drei Kammern, Kette, Blut & $5$ & $5$ & endlich\\
drei Kammern, voll & $9$ & $5$ & Dimension $4$\\
\bottomrule
\end{tabular}}
\end{center}

\begin{center}
\begin{minipage}{0.9\textwidth}
\textbf{Die Dimension der Faser ist die Zahl der Parameter, die die Daten nicht bestimmen. Sie ist die Differenz zweier Dimensionen --- und die des Bildes ist der Rang der Jacobi-Matrix der Abbildung, Kapitel 8.}
\end{minipage}
\end{center}

\begin{center}
\begin{minipage}{0.88\textwidth}
\kursiv{\textbf{\enquote{Für $w$ in einer offenen Teilmenge} ist das Kleingedruckte}: Über speziellen Datenpunkten kann die Faser größer sein --- die Ausnahmemengen von Kapitel 6. Für den Leitfall mit $k_{20} = 0$ ist die allgemeine Faser ein Punkt; aber wenn die gemessene Kurve nur eine Exponentialfunktion zeigt ($s_3 = 0$, also $k_{10}k_{21} = 0$), ist die Faser eine Gerade: Entweder $k_{10} = 0$ und $k_{21}$ beliebig mit $k_{12} = s_1$, oder umgekehrt. Ein Patient, dessen Kurve zufällig so aussieht, hat ein nicht identifizierbares Modell, obwohl das Modell \enquote{generisch identifizierbar} ist. Satz 7.5 sagt, dass solche Patienten selten sind; er sagt nicht, dass es sie nicht gibt.}
\end{minipage}
\end{center}
\end{merksatz}

\section{Der Bruchfall}

\begin{warnung}
\boxtitelr{Dimension ist keine Zahl, wenn die Varietät zerfällt}
$V(xz, yz) = $ Ebene $\cup$ Gerade hat Hilbert-Polynom $\tfrac12d^{2} + \tfrac52d + 1$: Dimension $2$, Grad $1$. Die Gerade ist verschwunden --- sie trägt nur zum linearen Term bei.

\begin{center}
{\small
\renewcommand{\arraystretch}{1.5}
\begin{tabular}{@{}P{3.2cm}P{5.6cm}@{}}
\toprule
\textbf{Dimension} & die größte Komponentendimension; Komponenten kleinerer Dimension sind unsichtbar\\
\textbf{Grad} & zählt nur die Komponenten größter Dimension\\
\textbf{Dimension in einem Punkt} & die Dimension der Komponenten durch den Punkt: $2$ auf der Ebene, $1$ auf der Geraden außerhalb, $2$ im Ursprung\\
\textbf{Der Leitfall} & Faser über $s_3 = 0$: zwei Geraden --- Dimension $1$, Grad $2$; die zwei Punkte der allgemeinen Faser sind zu zwei Geraden geworden\\
\bottomrule
\end{tabular}}
\end{center}

\begin{center}
\begin{minipage}{0.9\textwidth}
\textbf{Dimension und Grad sind für irreduzible Varietäten gemacht. Für reduzible sagen sie etwas über die größten Teile und nichts über die kleineren. Zuerst zerlegen, dann messen.}
\end{minipage}
\end{center}

\begin{center}
\begin{minipage}{0.88\textwidth}
\kursiv{\textbf{Varietäten, deren Komponenten alle dieselbe Dimension haben, heißen rein- oder äquidimensional}, und die meisten Sätze der Geometrie setzen das voraus, oft stillschweigend. Der Satz von Bézout (Kapitel 9) ist ein Beispiel: Er zählt Schnittpunkte von Kurven, und wenn eine der \enquote{Kurven} einen isolierten Punkt enthält (der Bruchfall von Kapitel 1 über $\mathbb R$: $y^{2} = x^{2}(x - 1)$), zählt er ihn mit --- über $\mathbb C$ ist der Punkt Teil der Kurve, und alles stimmt. Für den Leitfall: Die Faser mit zwei Geraden hat einen Schnittpunkt ($k_{10} = k_{21} = 0$), in dem die lokale Dimension $1$, aber der Tangentialraum (Kapitel 8) zweidimensional ist: die erste Singularität des Bandes.}
\end{minipage}
\end{center}
\end{warnung}

\section{Übungsaufgaben}

\begin{aufgabe}
\boxtitel{Aufgabe 7.1 --- Die verdrehte Kubik gezählt}
Zählen Sie die Standardmonome von $\mathrm{LM}(I) = \langle x^{2}, xy, xz, y^{3}\rangle$ (die Leitmonome der grevlex-Basis der verdrehten Kubik; grevlex liefert hier dieselben Leitmonome wie lex) und bestimmen Sie Dimension und Grad.

\loesung
\schritt{Schritt 1 --- Welche Monome sind Standard}
Ein Monom $x^{a}y^{b}z^{c}$ ist durch keines der vier teilbar, wenn: $a\le1$; ist $a = 1$, dann $b = c = 0$; und $b\le2$. Also: $x$ (allein) und $y^{b}z^{c}$ mit $b\in\{0, 1, 2\}$, $c\ge0$.

\schritt{Schritt 2 --- Zählen bis Grad $d$}
$y^{b}z^{c}$ mit $b + c\le d$: für $b = 0$: $d + 1$ Stück ($c = 0..d$); $b = 1$: $d$; $b = 2$: $d - 1$. Zusammen $3d$. Dazu $x$ (Grad $1\le d$): $3d + 1$ für $d\ge2$. Probe mit der Tabelle: $d = 3$: $10$. \checkmark

\schritt{Schritt 3 --- Dimension und Grad}
${}^{a}\!HP(d) = 3d + 1$: Grad $1$, Dimension $1$; Leitkoeffizient $3$, mal $1! = 3$: Grad $3$. Die verdrehte Kubik ist eine Kurve vom Grad $3$ --- daher der Name.

\schritt{Schritt 4 --- Geometrische Probe}
Eine allgemeine Ebene $ax + by + cz = e$ trifft $(t, t^{2}, t^{3})$ in den Nullstellen von $ct^{3} + bt^{2} + at - e$: drei Punkte. \checkmark

\begin{center}
\begin{minipage}{0.88\textwidth}
\kursiv{\textbf{Die Standardmonome $1, x, y, y^{2}, z, yz, y^{2}z, z^{2}, \ldots$ sind eine Basis des Koordinatenrings}, und ihre Struktur --- drei \enquote{Stränge} $y^{b}z^{c}$ und ein einsamer Zusatz $x$ --- ist die Dreiheit des Grades. Das Hilbert-Polynom einer Kurve vom Grad $e$ ist $ed + 1 - p$ mit einer Zahl $p$, dem Geschlecht; hier $p = 0$: Die verdrehte Kubik ist rational (Parametrisierung durch $t$), Kapitel 12.}
\end{minipage}
\end{center}
\end{aufgabe}

\begin{aufgabe}
\boxtitel{Aufgabe 7.2 --- Dimension des Durchschnitts}
Zwei Flächen im $\mathbb C^{3}$ schneiden sich in einer Kurve --- immer? Prüfen Sie an $V(z)\cap V(z - 1)$, $V(z)\cap V(z)$ und $V(xz, yz)\cap V(z - 1)$, und formulieren Sie die richtige Regel.

\loesung
\schritt{Schritt 1}
$V(z)\cap V(z - 1) = \emptyset$: zwei parallele Ebenen. Dimension \enquote{$-1$}.

\schritt{Schritt 2}
$V(z)\cap V(z) = V(z)$: dieselbe Ebene, Dimension $2$.

\schritt{Schritt 3}
$V(xz, yz)\cap V(z - 1)$: auf $z = 1$ ist $x = y = 0$: ein Punkt, Dimension $0$. Die Ebene $z = 0$ wird gar nicht getroffen, die Gerade in einem Punkt.

\schritt{Schritt 4 --- Die Regel}
Für irreduzible $V, W\subseteq\mathbb C^{n}$ hat jede Komponente von $V\cap W$ Dimension $\ge\dim V + \dim W - n$ --- falls der Durchschnitt nicht leer ist (affiner Dimensionssatz, Auslassung 8). Schritt 1: leer, die Regel sagt nichts. Schritt 2: $2 + 2 - 3 = 1\le2$. \checkmark\ Schritt 3: $V$ reduzibel --- für die Gerade: $1 + 2 - 3 = 0\le0$. \checkmark\ Für die Ebene: Durchschnitt leer.

\begin{center}
\begin{minipage}{0.88\textwidth}
\kursiv{\textbf{\enquote{Falls nicht leer} ist der Haken, den die projektive Geometrie entfernt}: Im projektiven Raum schneiden sich zwei Flächen immer (Kapitel 9), und die Ungleichung gilt ohne Ausnahme. Für den Leitfall ist der Dimensionssatz das Zählargument aus Aufgabe 6.3 in Reinform: Drei Gleichungen im $\mathbb C^{4}$, jede eine Hyperfläche der Dimension $3$, schneiden sich in Dimension $\ge 3 + 3 + 3 - 2\cdot4 = 1$ --- eine Kurve, mindestens.}
\end{minipage}
\end{center}
\end{aufgabe}

\begin{aufgabe}
\boxtitel{Aufgabe 7.3 --- Drei Kammern}
Ein Drei-Kammer-Modell in Kette (Blut $\leftrightarrow$ Gewebe $1\leftrightarrow$ Gewebe $2$, Ausscheidung nur aus dem Blut) hat die Parameter $k_{10}, k_{12}, k_{21}, k_{23}, k_{32}$. Die Blutkurve ist eine Summe dreier Exponentialfunktionen. Wie viele Messgrößen, wie groß die Faser, und was sagt das Zählen nicht?

\loesung
\schritt{Schritt 1 --- Messgrößen}
Drei Exponenten, drei Amplituden, eine Amplitudensumme durch die Dosis festgelegt: $6 - 1 = 5$ Messgrößen.

\schritt{Schritt 2 --- Faser}
$\varphi\colon\mathbb C^{5}\to\mathbb C^{5}$. Ist das Bild dicht --- keine polynomiale Beziehung zwischen den fünf Messgrößen ---, so hat die allgemeine Faser Dimension $0$: endlich viele Modelle, das Modell ist (lokal) identifizierbar.

\schritt{Schritt 3 --- Ob das Bild dicht ist}
Das sagt das Zählen nicht. Man muss nachrechnen, dass die fünf Messgrößen als Funktionen der fünf Parameter unabhängig sind --- Rang der Jacobi-Matrix $5$ (Kapitel 8) oder Eliminationsideal $\{0\}$ (Kapitel 6). Für die Kette ist es so (das Modell ist ein Standardfall der Pharmakokinetik und global identifizierbar); für einen Ring aus drei Kammern mit sechs Übergängen wäre $\dim V = 7 > 5$ und die Faser mindestens zweidimensional.

\schritt{Schritt 4 --- Wie viele Punkte}
Auch das sagt das Zählen nicht; das sagt der Grad der Faser, also die Hilbert-Funktion des Ideals der Messgleichungen --- und für die Kette ist er $1$: global identifizierbar. Für die Variante mit Ausscheidung aus Gewebe $2$ statt aus dem Blut ist er größer als $1$.

\begin{center}
\begin{minipage}{0.88\textwidth}
\kursiv{\textbf{Zählen gibt eine Schranke, Dimension gibt die Antwort auf \enquote{ob}, Grad die Antwort auf \enquote{wie viele}} --- drei Stufen, und jede braucht ein eigenes Werkzeug. Die Programme der strukturellen Identifizierbarkeit durchlaufen genau diese drei Stufen; die Pharmakologie der siebziger Jahre blieb auf der ersten stehen und hielt manches Modell für bestimmt, dessen Faser zwei Punkte hatte. Das Zwei-Kammer-Modell mit Messung im Gewebe (Aufgabe 1.3) ist ein solcher Fall, und er ist in Lehrbüchern bis heute nicht immer korrekt wiedergegeben.}
\end{minipage}
\end{center}
\end{aufgabe}

\begin{geschichte}
\boxtitelg{Hilbert 1890, nochmals --- und Macaulay}
Das Hilbert-Polynom steht in derselben Arbeit von 1890 wie der Basissatz: Hilbert zeigte, dass die Zahl der linear unabhängigen Formen vom Grad $d$ modulo einem Ideal für große $d$ ein Polynom ist, mit Hilfe der Syzygien --- der Beziehungen zwischen den Erzeugern --- und ihrer Syzygien, bis die Kette nach höchstens $n$ Schritten abbricht (der Syzygiensatz).

\begin{center}
\begin{minipage}{0.9\textwidth}
\kursiv{\textbf{\Person{Francis Macaulay}, Gymnasiallehrer in London, bewies 1927 den Satz 7.2}: Jedes Ideal hat dieselbe Hilbert-Funktion wie sein Leitmonomideal --- und charakterisierte, welche Funktionen überhaupt als Hilbert-Funktionen vorkommen (eine Bedingung an die Binomialdarstellung der Werte, die heute Macaulays Satz heißt). Seine Arbeit wurde erst in den sechziger Jahren gelesen; \Person{Stanley} benutzte sie 1975 für die Kombinatorik konvexer Polytope, und das Computeralgebrasystem für kommutative Algebra, das \Person{Grayson} und \Person{Stillman} 1993 schrieben, heißt zu seinen Ehren Macaulay2.}
\end{minipage}
\end{center}

\schritt{Krull}
\begin{center}
\begin{minipage}{0.9\textwidth}
\kursiv{\Person{Wolfgang Krull} definierte 1928 die Dimension eines Rings über Primidealketten --- eine Definition ohne Geometrie, die für Zahlringe ebenso gilt wie für Koordinatenringe ($\mathbb Z$ hat Dimension $1$). \textbf{Dass sie mit Hilberts Polynomgrad und dem Transzendenzgrad übereinstimmt, ist der Hauptsatz der Dimensionstheorie}; sein Beweis über Noethers Normalisierungslemma von 1926 ist Auslassung 6 und steht in jedem Lehrbuch der kommutativen Algebra, am kürzesten bei \Person{Atiyah} und \Person{Macdonald} 1969, Kapitel 11.}
\end{minipage}
\end{center}
\end{geschichte}

\begin{forschung}
\boxtitelp{Dimension heute}
\begin{center}
{\small
\renewcommand{\arraystretch}{1.45}
\begin{tabular}{@{}P{3.0cm}P{5.2cm}@{}}
\toprule
\textbf{Numerische Dimension} & Dimension einer Varietät aus zufälligen Schnitten mit linearen Unterräumen bestimmen, numerisch (Sommese--Wampler 2005): ohne Gröbner-Basis, für Systeme mit hundert Veränderlichen.\\
\textbf{Algebraische Statistik} & Die Dimension eines statistischen Modells (als Varietät im Wahrscheinlichkeitssimplex) ist die Zahl seiner wirklich freien Parameter; Modelle mit versteckten Variablen haben oft kleinere Dimension als Parameter --- Nichtidentifizierbarkeit, Kapitel 13.\\
\textbf{Intrinsische Dimension von Daten} & Band 28, Kapitel 11: Dieselbe Frage für Punktwolken statt Varietäten, mit Schätzern statt Hilbert-Polynomen; die algebraische Fassung ist der Grenzfall ohne Rauschen.\\
\bottomrule
\end{tabular}}
\end{center}

\schritt{Zur zweiten Zeile}
\begin{center}
\begin{minipage}{0.88\textwidth}
\kursiv{\textbf{Ein Modell mit versteckten Variablen --- Mischungen, Faktoren, latente Klassen --- ist eine Abbildung von Parametern in Wahrscheinlichkeiten, und die Dimension seines Bildes ist, was die Daten je sehen können.} Für Mischungen von zwei Binomialverteilungen, für Faktormodelle mit einem Faktor, für phylogenetische Bäume: Die Dimension des Bildes ist berechnet, und wo sie unter der Parameterzahl liegt, ist das Modell nicht identifizierbar, egal wie viele Daten man hat. Band 17 hat diese Fälle als Warnungen genannt; Band 29 sagt, dass sie eine Dimension sind, und wie man sie ausrechnet. Der Leitfall ist das einfachste Beispiel der Gattung.}
\end{minipage}
\end{center}
\end{forschung}

\begin{merksatz}
\boxtitel{Was Sie jetzt können}
\begin{center}
\renewcommand{\arraystretch}{1.5}
\begin{tabular}{@{}P{3.0cm}P{5.6cm}@{}}
\toprule
\textbf{Definition 7.1} & Hilbert-Funktion ${}^{a}\!H_I(d)$\\
\textbf{Satz 7.2} & Macaulay: Standardmonome zählen (gradkompatibel)\\
\textbf{Satz 7.3} & Hilbert-Polynom; $\dim = $ Grad, $\deg = \dim!\cdot$Leitkoeffizient\\
\textbf{Satz 7.4} & Hilbert = Krull = Transzendenz\\
\textbf{Satz 7.5} & Faserdimension $\ge\dim V - \dim W$, generisch $=$\\
\textbf{Der Bruchfall} & reduzibel: Dimension sieht nur die größten Teile\\
\bottomrule
\end{tabular}
\end{center}

\begin{center}
\begin{minipage}{0.88\textwidth}
\kursiv{Wie groß --- beantwortet. Kapitel 8 fragt, wie glatt: Wo hat eine Varietät eine Tangente, wo eine Spitze, wo kreuzt sie sich --- und was die Jacobi-Matrix damit zu tun hat, mit der die Identifizierbarkeit in der Praxis geprüft wird.}
\end{minipage}
\end{center}
\end{merksatz}
